\section{Ideals and Quotients: Remarks and Examples. Prime and Maximal Ideals}
\subsection{Basic Operations}
We often define ideals in terms of a set of generators. Let $a\in R$ be any element. Then the subsets $Ra$ and $aR$ are left- and right-ideals, respectively. If the ring is commutative, these sets are clearly the same, and we denote them $(a)$.
\begin{definition}[Principal Ideal]\label{dfn:3.4.1}
	Let $R$ be a ring and let $a\in R$. Then the principal ideal generated by $a$ is the set
	\[
		(a)=aR=Ra
	.\]
\end{definition}

From the exercises, we know that for any collection $\left\{ I_\alpha  \right\}_{\alpha \in A}$ of ideals of $R$, the sum
\[
	\sum_{\alpha\in A } I_\alpha 
\]
is an ideal of $R$. We have a similar idea for ideals generated by multiple elements. Let $\left\{ a_{\alpha } \right\}_{\alpha \in A}\subseteq R$. Then define
\[
	\left( a_{\alpha } \right)_{\alpha \in A}\coloneqq \sum_{\alpha \in A} \left( a_{\alpha } \right) 
.\]
In the finite case, this is given as
\[
	\left( a_1, \dots, a_n  \right) =(a_1)+\cdots +(a_n)
.\]
This is the smallest ideal containing $\left\{ a_{\alpha } \right\}_{\alpha \in A}$. An ideal $I$ is called \emph{finitely generated} if there exists a finite collection $a_1, \dots, a_n \in R$ such that
\[
	I=\left( a_1, \dots, a_n  \right) 
.\]

One example that I want to copy down is as follows. Let $R$ be a commutative ring, and let $a,b\in R$. Let $\overline{b}  $ denote the equivalence class of $b$ under the relation $R /(a)$. Then
\[
	\frac{R /(a)}{(\overline{b} )}\cong \frac{R}{(a,b)}
.\]
Note that this follows from proposition \ref{prp:3.3.2}, as for some reason,
\[
	(\overline{b} )\cong \frac{(a,b)}{(b)}
.\]
I don't really care to prove this honestly.

Principal ideals are very special, so we use special terminology when referring to them.
\begin{definition}[Noetherian]\label{dfn:3.3.2}
	A commutative ring $R$ is \emph{Noetherian} if every ideal of $R$ is finitely generated.
\end{definition}
\begin{definition}[Principal Ideal Domain]\label{dfn:3.4.3}
	An integral domain $R$ is a PID (Principal Ideal Domain) if every ideal of $R$ is principal.
\end{definition}
\begin{proposition}\label{prp:3.4.1}
	$\mathbb{Z}$ is a PID.
\end{proposition}
\begin{proof}
	Let $I\subseteq \mathbb{Z}$ be an ideal. Since $I$ is a subgroup, $I=n\mathbb{Z}=(n)$ for some $n\in\mathbb{Z}$.
\end{proof}

Interestingly, the fact that $\mathbb{Z}$ is a PID explains why gcds behave as they do. If $m,n\in\mathbb{Z}$, then $(m,n)$ is principal, meaning there exists $d\in \mathbb{Z}$ such that $(m,n)=(d)$. This $d$ is fairly trivially $\operatorname{gcd}(m,n) $. Additionally, for any field $k$, the ring $k[x]$ is also a PID, which we will prove soon. By contrast, $\mathbb{Z}[x]$ is not a PID. For example, $(2,x)$ cannot be generated by a single element.

Here are two more operations on ideals. Let $\left\{ I_{\alpha } \right\}_{\alpha \in A}$ be a collection of ideals of some ring $R$. Then
\[
	\bigcap_{\alpha \in A} I_{\alpha }
\]
is an ideal of $R$; it is the largest ideal contained in $\left\{ I_{\alpha } \right\}_{\alpha \in A}$. Another operation is the following: $IJ$ denotes the ideal generated by all products $ij$ with $i\in I$, $j\in J$. Note that this is not the same as the $IJ$ notation used in group theory, as it is not just the products in these ideals, but rather the ideal it generates. Moreover, note that $ij\in I,J$ since $I,J$ are ideals, so we have $IJ\subseteq I\cap J$.
\subsection{Quotients of Polynomial Rings}
Let $R$ be a ring, and let $f(x) = x^d + \cdots + a_1x^1 + a_0\in R[x]$. We are assuming for convenience that $f$ is monic; the leading coefficient is the identity $1_R$. It will be shown in the exercises that  if $R$ is a field, then all principal ideals are generated by unique monic polynomials. In the general case, however, choosing $f$ to be monic can have impacts on the ideal. However, assuming it is monic is convenient, because it is necessarily not a zero divisor, and additionally we have a notion of division in general rings. That is, for all $g(x)\in R[x]$, there exist $q(x),r(x)\in R[x]$ such that
\[
	g(x)=q(x)p(x)+r(x)
\]
and $\operatorname{deg}r < \operatorname{deg}f  $, where we take the convention that the polynomial $0$ has degree $-\infty$.
\begin{lemma}\label{lma:3.4.1}
	Let $f(x)$ be a monic polynomial, and suppose
	\[
		f(x)q_1(x)+r_1(x)=f(x)q_2(x)+r_2(x)
	,\]
	where $\operatorname{deg}r_1,\operatorname{deg}r_2 < \operatorname{deg}f   $. Then $q_1(x)=q_2(x)$ and $r_1(x)=r_2(x)$.
\end{lemma}
\begin{proof}
	We have
	\[
		f(x)(q_1(x)-q_2(x))=r_1(x)-r_2(x)
	.\]
	Suppose $r_2 \neq r_1$. Then $\operatorname{deg}(r_2-r_1) <\operatorname{deg}f $. However, since $f(x)$ is monic,
	\[
		\operatorname{deg}(f(x)q(x)) =\operatorname{deg} f(x)+\operatorname{deg} q(x)
	.\]
	Thus, $\operatorname{deg}(f(q_1-q_2)) > \operatorname{deg} f$, which is a contradiction. Therefore, $r_1(x)=r_2(x)$. Since $f$ is monic and thus not a zero divisor, it follows that $q_1=q_2$.
\end{proof}

Note here that although division appears to be unique, it is only really unique in the commutative case. For noncommutative rings, we may have distinct left- and right-divisors of $f$. We will only consider the commutative case here, largely for convenience.

This is equivalent to the following: for all $g(x)\in R[x]$, there exists a unique $r(x)\in R[x]$ with $\operatorname{deg} r < \operatorname{deg} f$ such that
\[
	g(x)+(f(x))=r(x)+(f(x))
.\]

We can refine this observation. Note that any polynomial with degree less than $d$ can be viewed as an element of the direct sum $R ^{\oplus d}$. The map
\[
	\psi (r_1, \dots, r_{d-1}) = r_0 + r_1x + \cdots + r_{d-1}x ^{d-1}
.\]
is clearly a \emph{group} homomorphism. We also have that it is injective fairly trivially. We would like $\psi $ to be a \emph{ring} homomorphism, but the multiplicative structure of $R[x]$ and $R^d$ differ. We can reconsile this by the following proposition.
\begin{proposition}\label{prp:3.4.2}
	Let $R$ be a commutative ring, and let $f(x)\in R[x]$ have degree $d$. Then the map $\phi : R[x] \to R ^{\oplus d}$ defined by mapping $g(x)$ to its remainder when dividing by $f(x)$ induces an isomorphism of groups
	\[
		\frac{R[x]}{(f(x))}\cong R ^{\oplus d}
	.\]
\end{proposition}
\begin{proof}
	By lemma \ref{lma:3.4.1}, $\phi $ is well-defined. We also have that $\phi $ is surjective, because $\psi $ as previously defined is a right-inverse:
	\[
		(\phi \circ \psi  )(r_0, \dots, r_{d-1})=\phi (r_0+r_1x + \cdots + r_{d-1}x ^{d-1}) = (r_0, \dots, r_{d-1})
	,\]
	as $f$ cannot divide into $r(x)$. By corollary \ref{cor:2.8.1},
	\[
		R ^{\oplus d}\cong \frac{R[x]}{\ker \phi }
	.\]
	Note that $\phi (g(x))=0$ if and only if the remainder is $0$; equivalently, $f(x)q(x)=g(x)$ for some $q(x)\in R[x]$, and thus $g(x)\in (f(x))$. Therefore, $\ker \phi =(f(x))$.
\end{proof}

This allows us to define multiplication in a new way on the direct sum $R ^{\oplus d}$, by using this isomorphism of groups and extending it in the most natural way to an isomorphism of rings. For a simple example, let $f(x)=x-a$ have degree 1, and note that $r$ must thus be a constant. Then note that $\phi $ is simply the evaluation map at $a$. We have
\[
	g(x)=(x-a)q(x) + r \implies g(a)=(a-a)q(x) + r = r
.\]
This induces an isomorphism
\[
	\frac{R[x]}{(x-a)}\cong R
,\]
which is in fact an isomorphism of rings. We can extend this idea to more complex examples, as for every monic polynomial, there is a brand new ring structure on $R ^{\oplus d}$. For $d=1$ these structures are all the same, but new structures arise as soon as $d=2$. Here is a concrete example.

Let $f(x)=x^2 + 1_R$. We have by proposition \ref{prp:3.4.2} that
\[
	R \oplus R \cong \frac{R[x]}{(x^2 + 1_R)}
.\]
We can analyze the multiplication operation induced on $R\oplus R$. Recall that each element of the ideal can be represented as a single remainder polynomial; in this case, that is a polynomial of degree one, of the form $r(x)=a_0 + a_1x$. We have
\[
	(a_0+a_1x)(b_0+b_1x)=a_0b_0 + (a_0b_1+a_1b_0)x+a_1b_1x^2
.\]
Dividing by $(x^2+1_R)$, we have
\[
	a_0b_0+(a_0b_1+a_1b_0)x+a_1b_1x^2 = \left( x^2+1 \right)a_1b_1 + ((a_0b_0-a_1b_1)+(a_0b_1+a_1b_0)x)
.\]
Applying $\phi $, we have
\[
	\phi (\left( a_0+a_1x \right) \left( b_0+b_1x \right) )=\left( a_0b_0-a_1b_1,a_0b_1+a_1b_0 \right) 
.\]
This induces the multiplication structure
\[
	(a_0,a_1)(b_0,b_1)=(a_0b_0-a_1b_1,a_0b_1+a_1b_0)
\]
on $R\oplus R$. One might recognize this as precisely the multiplication operation in $\mathbb{C}$, where an element of $\mathbb{C}$ is identified with a pair of real numbers $(a_0,a_1)\in \mathbb{R}\oplus \mathbb{R}$. That is,
\[
	\frac{\mathbb{R}[x]}{(x^2+1)}\cong \mathbb{C}
\]
as rings. This is not suprising, as $x^2+1$ does not have real roots, and since the roots will show up as classes of $x$ and $-x$ in the quotient, the construction should contain the roots. We will revisit these constructions to solve equations in chapters 5 and 7.

\subsection{Prime and Maximal Ideals}
We are still assuming that rings are commutative, since we will not use these definitions in noncommutative contexts.
\begin{definition}[Prime/Maximal Ideal]\label{dfn:3.4.4}
	Let $I\neq (1)$ be an ideal of a commutative ring $R$. $I$ is called a prime ideal if $R /I$ is an integral domain, and $I$ is called a maximal ideal if $R /I$ is a field.
\end{definition}

These definitions are equivalent to the following proposition, and it is a matter of preference whether to define these in terms of quotients or this proposition.
\begin{proposition}\label{prp:3.4.3}
	Let $I\neq (1)$ be an ideal of a commutative ring $R$. Then
	\begin{itemize}
		\item $I$ is prime if and only if for all $a,b\in R$,
			\[
				ab\in I \implies (a\in I \lor b\in I)
			.\]
		\item $I$ is maximal if and only if for all ideals $J$ of $R$,
			\[
				I\subseteq J\implies (I=J \lor J=R)
			.\]
	\end{itemize}
\end{proposition}
\begin{proof}
	The ring $R$ is an integral domain if and only if for all $a+I,b+I\in R /I$,
	\[
		ab+I=I \implies (a\in I \lor b\in I)
	.\]
	Wow, this is an equivalent statement.

	From the exercises, we know a commutative ring is a field if and only if its only ideals are $\left( 0 \right) $ and $(1)$. Using this and the fact that there is a bijection between ideals containing $I$ and ideals of $R /I$ (proposition \ref{prp:3.3.2}), we have that $R$ is a field.
\end{proof}

Note that all maximal ideals are prime, but the converse is not necessarily true. There is, however, a case where it holds.
\begin{proposition}\label{prp:3.4.4}
	Let $R$ be a ring with an ideal $I$. If $R /I$ is finite, then $I$ is prime if and only if $I$ is maximal.
\end{proposition}
\begin{proof}
	Follows immediately from proposition \ref{prp:3.1.3}.
\end{proof}

One example is that $\mathbb{Z}/n\mathbb{Z}$ is finite, so $(n)$ is prime/maximal when $n$ is prime. Attributing this uniqueness to $\mathbb{Z}$ is mildly misleading, however, as it holds for all PIDs.
\begin{proposition}\label{prp:3.4.5}
	Let $R$ be a PID,  and let $I\neq (0)$ be an ideal of $R$. Then $R$ is prime if and only if $R$ is maximal.
\end{proposition}
\begin{proof}
	Maximal ideals are always prime, so we show the converse. Let $I=(a)$ be prime, and suppose $I\subseteq J$ is an ideal of $R$. Note that $J$ must be a principal ideal, so  $J=(b)$.  Since $(a)\subseteq (b)$, we have that $a\in (b)$, and thus $a=bc$ for some $c\in R$. But then either $b\in (a)$ or $c\in (a)$ since $(a)$ is prime. If $b\in (a)$, then $(b)\subseteq (a)$. If $c\in (a)$, then $c=da$ for some $d\in R$, so $a=bc=bda$, and thus  $bd=1$. Therefore, $b$ is a unit, so $1\in (b)=R$.
\end{proof}

\begin{definition}[Spectrum]\label{dfn:3.4.5}
	The spectrum of a commutative ring $R$, denoted $\spec R$, is the set of all prime ideals in $R$.
\end{definition}

The commentary on the spectrum involves figures too advanced for me to bother copying, so I recommend reading the tailend of section III.4.3 in the book. We can draw the spectrum, where each level higher contains the lower levels I believe. Drawing the spectrum of $k[x]$ for $k$ a field would look a lot like $\mathbb{Z}$; all the prime numbers at one level, all containing the nonmaximal ideal $(0)$.

This picture is nice for algebraicly closed fields, like $\mathbb{C}$. Using this fact, it is shown in the exercises that the maximal ideals in $\mathbb{C}[x]$ are all the ideals of the form $(x-z)$. Drawing $\spec\mathbb{C}$, we draw one level of prime/maximal ideals above the ideal $\left( 0 \right) $, so we say $\mathbb{C}[x]$ has Krull dimension 1.

\begin{definition}[Krull Dimension]\label{dfn:3.3.6}
	The Krull dimension of a ring $R$ is the longest chain of inclusion of prime ideals in $R$.
\end{definition}

For example, any field $k$ has Krull dimension $0$. More generally, we can show that $k[x_1, \dots, x_n ]$ has Krull dimension $n$. The chain of length $n$ is the chain
\[
	(0) \subsetneq (x_1)\subsetneq (x_1,x_2)\subsetneq \cdots \subsetneq (x_1, \dots, x_n )
.\]
These ideals are shown to be prime in the exercises.

In the future, we will explore these ideals in more detail. It can be shown that the maximal ideals of $\mathbb{C}[x_1, \dots, x_n ]$ are all of the ideals of the form $(x_1-z_1, \ldots, x_n-z_n)$, but this requires a deep result known as \emph{Hilbert's Nullstellensatz}. We will return to this in chapter 7.
